\begin{helper}[Advance a repeated missing color along the path suffix]
Let \((\varphi,T)\) be a minimum bad pair with
\(T=(y_0,e_1,y_1,\ldots,e_n,y_n)\).  Suppose \(i+1<n\), the color
\(\beta\) is missing at both \(y_i\) and \(y_n\), and either \(i=0\)
or \(\rho(T)\leq i\).  Then there is a coloring \(\varphi'\) such that
\((\varphi',T)\) is again a minimum bad pair and \(\beta\) is missing
at both \(y_{i+1}\) and \(y_n\) under \(\varphi'\).
\end{helper}
\begin{proof}
Unfolding \noderef{TashTreeMinimumBadPair} and
\noderef{TashTreeBadPair}, and then \noderef{KTriple}, supplies
properness of \(\varphi\), the bound
\(\Delta(G)+1\leq k\), and the Tashkinov condition for \(T\).
By \noderef{MissingColorsNonemptyOfMaximumDegree}, applied with that
degree bound at \(y_{i+1}\), choose a color \(\gamma\) missing there.

The colors \(\beta,\gamma\) are distinct.  Otherwise both \(y_i\) and
\(y_{i+1}\) would miss \(\beta\).  These are distinct vertices by
injectivity in \noderef{TreeSequence}, and both lie in the proper prefix
through \(y_{i+1}\) because \(i+1<n\).  The elementariness supplied for
that prefix by \noderef{TashTreeProperPrefixElementary}, together with
the defining disjointness of missing-color sets in
\noderef{ColoredVertexSetElementary}, would be contradicted.

We next prove the two \noderef{TreeColorUnusedThrough} hypotheses
through \(y_{i+1}\), constructing the relevant exact prefix with
\noderef{TreeSequencePrefix}.  If \(i=0\), the formal range in either
unused-through assertion consists of edge indices \(r\) with
\(1\leq r<i+1=1\), so it is empty.  Both assertions are therefore
vacuous in this boundary case, and no endpoint equation for the initial
edge \(e_1\) is used.

Suppose now that \(i>0\).  If \(\gamma\) occurred on one of the
displayed colored edges \(e_2,\ldots,e_{i+1}\), the
\noderef{TashkinovTree} condition would provide an earlier displayed
vertex missing \(\gamma\).  That vertex is distinct from \(y_{i+1}\),
so the elementary proper prefix through \(y_{i+1}\) would contain two
vertices missing \(\gamma\), contrary to
\noderef{ColoredVertexSetElementary}.  If \(\beta\) occurred on one of
\(e_2,\ldots,e_i\), its earlier Tashkinov witness would similarly be a
vertex distinct from \(y_i\) in the same proper prefix.  Elementariness
then contradicts the fact that \(y_i\) also misses \(\beta\).

It remains to exclude \(\beta\) from \(e_{i+1}\).  By the
least-index specification in \noderef{TashTreePathIndex}, the suffix
starting at \(\rho(T)\) is a path.  Our hypothesis \(\rho(T)\leq i\)
therefore lets us apply \noderef{TashTreeSuffixPath} at the zero-based
edge index \(i\): the exact edge \(e_{i+1}\) joins \(y_i\) to
\(y_{i+1}\).  This edge cannot have color \(\beta\), since it is
incident with \(y_i\), where \(\beta\) is missing by
\noderef{PresentColors} and \noderef{MissingColors}.  Thus both colors
are unused through \(y_{i+1}\).

Apply \noderef{TashTreeControlledInterchange} to the distinct colors
\(\beta,\gamma\), the vertices \(y_i,y_{i+1}\), and these unused-color
facts.  It returns the exact component-swap coloring \(\varphi'\), says
that \((\varphi',T)\) is a minimum bad pair, and makes \(\beta\) missing
at \(y_{i+1}\).  It remains to verify that \(\beta\) is still missing
at \(y_n\).

The controlled interchange says that the swapped component is a path
component having \(y_i,y_{i+1}\) as its two vertex ends.  The vertex
\(y_n\) is distinct from both by \(i+1<n\) and vertex injectivity in
\noderef{TreeSequence}.  Since \(y_n\) misses \(\beta\),
\noderef{MissingColorIncidenceEndpoint} gives its vertex-node degree at
most one in the same two-color incidence graph.  If \(y_n\) lay in the
swapped component, the simple path supplied by
\noderef{TwoColorComponentSpanningPath} would consequently have the
three distinct ends \(y_i,y_{i+1},y_n\), which is impossible.  Hence
\(y_n\) is outside the swapped component.  The outside-vertex clause of
\noderef{TwoColorComponentSwapMissingBehavior} fixes its missing-color
set, so \(\beta\) remains missing there under \(\varphi'\).  This proves
the claim.
\end{proof}
